\chapter{Fractions inside fractions}
\chaptertag{Continued fractions and best approximations}

The quotients from Euclid's algorithm are not bookkeeping debris. Read them forward and they produce a new numerical language.

\section{Euclid, folded up}

From
\[
415=4\cdot93+43,\quad 93=2\cdot43+7,\quad
43=6\cdot7+1,\quad 7=7\cdot1,
\]
we obtain
\[
\frac{415}{93}=4+\cfrac1{2+\cfrac1{6+\cfrac1{7}}}=[4;2,6,7].
\]

\begin{definition}
A finite simple continued fraction is
\[
[a_0;a_1,\dots,a_n]
=a_0+\cfrac1{a_1+\cfrac1{\ddots+\cfrac1{a_n}}},
\]
where $a_0\in\Z$ and $a_i\in\N$ for $i\ge1$.
\end{definition}

The truncations $[a_0;\dots,a_k]=p_k/q_k$ are called \emph{convergents}. They can be computed without rebuilding the nested fraction:
\[
p_k=a_kp_{k-1}+p_{k-2},\qquad q_k=a_kq_{k-1}+q_{k-2},
\]
with $p_{-2}=0,p_{-1}=1,q_{-2}=1,q_{-1}=0$.

\section{Linear fractional transformations}
\label{sec:mobius}

A \emph{linear fractional transformation}, also called a M\"obius
transformation, is a function of the form
\[
 T_M(t)=\frac{at+b}{ct+d},
 \qquad M=\begin{pmatrix}a&b\\c&d\end{pmatrix},
 \qquad ad-bc\ne0.
\]
We will use integer or rational coefficients. The formula is defined when
$ct+d\ne0$. The determinant condition excludes a constant function and
ensures that the transformation can be undone.

Why does a $2\times2$ matrix act on a single number? Represent $t$ by the
column $(t,1)^T$, multiply, and take the ratio of the two coordinates:
\[
 \begin{pmatrix}a&b\\c&d\end{pmatrix}
 \begin{pmatrix}t\\1\end{pmatrix}
 =\begin{pmatrix}at+b\\ct+d\end{pmatrix}
 \quad\longmapsto\quad \frac{at+b}{ct+d}.
\]
Thus multiplying $M$ by any nonzero scalar does not change $T_M$. In this
chapter, $T_M(t)$ means this ratio, not ordinary multiplication of a matrix
by a scalar. One may also include a point $\infty$: when $c\ne0$,
$T_M(\infty)=a/c$ and $T_M(-d/c)=\infty$; when $c=0$, $\infty$ stays fixed.

\begin{proposition}[Composition is matrix multiplication]
For invertible $2\times2$ matrices $M,N$,
\[
 T_M\circ T_N=T_{MN},\qquad (T_M)^{-1}=T_{M^{-1}}.
\]
\end{proposition}
\begin{proof}
If $N=\begin{psmallmatrix}e&f\\g&h\end{psmallmatrix}$, substitution gives
\[
 T_M(T_N(t))
 =\frac{a(et+f)+b(gt+h)}{c(et+f)+d(gt+h)}
 =\frac{(ae+bg)t+(af+bh)}{(ce+dg)t+(cf+dh)}.
\]
These four coefficients are the entries of $MN$. Taking $N=M^{-1}$
proves the inverse rule. Explicitly, solving $y=(at+b)/(ct+d)$ gives
$t=(dy-b)/(a-cy)$.
\end{proof}

As with functions, the rightmost factor acts first. In particular,
\[
 D(a)=\begin{pmatrix}a&1\\1&0\end{pmatrix}
 \quad\text{acts by}\quad T_{D(a)}(t)=a+\frac1t.
\]
Therefore
\[
 [a_0;a_1,\ldots,a_k,t]=T_{D(a_0)D(a_1)\cdots D(a_k)}(t).
\]
Here the final entry $t$ is an entire remaining tail, not necessarily an
integer. The nesting of a continued fraction has become a product of
matrices.

\section{The matrix engine}

Package one continued-fraction digit into the matrix
\[
  D(a)=\begin{pmatrix}a&1\\1&0\end{pmatrix}.
\]
Repeated matrix multiplication stores every convergent at once:
\begin{equation}\label{eq:cf-matrix-product}
 D(a_0)D(a_1)\cdots D(a_k)
 =\begin{pmatrix}p_k&p_{k-1}\\q_k&q_{k-1}\end{pmatrix}.
\end{equation}
Indeed, multiplying the right-hand side by $D(a_{k+1})$ produces exactly the
two convergent recurrences. The first column advances to the next convergent,
while the old first column becomes the second.

Each digit matrix has determinant $-1$, not $+1$. Thus it lies in
$\operatorname{GL}_2(\Z)$, and a product with an even number of factors lies
in $\operatorname{SL}_2(\Z)$. More important than the sign is
\emph{unimodularity}: every product has determinant $\pm1$ and therefore has
an inverse with integer entries. Appendix~\ref{app:matrices} develops the
small amount of matrix algebra behind this observation.  Appendix~\ref{app:gosper}
then uses the same M\"obius machinery to perform exact arithmetic directly on
continued-fraction streams.

For example,
\[
 D(4)D(2)D(6)D(7)
 =\begin{pmatrix}415&58\\93&13\end{pmatrix},\qquad
 \begin{pmatrix}415&58\\93&13\end{pmatrix}^{-1}
 =\begin{pmatrix}13&-58\\-93&415\end{pmatrix}.
\]
The entries of the inverse already display the certificate
$1=13\cdot415-58\cdot93$. This is the first payoff of invertibility; the next
section explains why it is exactly the extended Euclidean algorithm.

\section{Extended Euclid in the convergents}

The convergents do not merely reuse Euclid's quotients: their numerators and
denominators are exactly the coefficients produced by the extended Euclidean
algorithm.

\begin{theorem}[Remainders from convergents]
Let $A>B>0$, and write the Euclidean divisions as
\[
 r_{-2}=A,\qquad r_{-1}=B,\qquad
 r_{k-2}=a_k r_{k-1}+r_k\quad(k=0,1,\dots,n),
\]
where $r_{n-1}=\gcd(A,B)$ and $r_n=0$. If
\[
  \frac AB=[a_0;a_1,\dots,a_n]
  \quad\text{and}\quad
  \frac{p_k}{q_k}=[a_0;a_1,\dots,a_k],
\]
then
\[
  \boxed{r_k=(-1)^k(q_kA-p_kB)}\qquad(0\le k\le n).
\]
In particular, the penultimate convergent gives a B\'ezout identity for the gcd.
\end{theorem}

\begin{proof}
For $k=0$, the formula reads $r_0=A-a_0B$, the first Euclidean division. For
$k=1$,
\[
  r_1=B-a_1r_0=-a_1A+(a_0a_1+1)B=-q_1A+p_1B.
\]
If the formula holds for the preceding two remainders, then
\begin{align*}
r_k&=r_{k-2}-a_kr_{k-1}\\
&=(-1)^k\bigl((a_kq_{k-1}+q_{k-2})A
 -(a_kp_{k-1}+p_{k-2})B\bigr)\\
&=(-1)^k(q_kA-p_kB),
\end{align*}
using the convergent recurrences. This proves the formula by induction.
\end{proof}

For $A/B=415/93$, the two algorithms can be read in the same table:
\begin{center}
\small
\begin{tabular}{cccc}
\toprule
$k$ & $a_k$ & convergent $p_k/q_k$ & signed linear combination $r_k$ \\
\midrule
$0$ & $4$ & $4/1$ & $415-4(93)=43$ \\
$1$ & $2$ & $9/2$ & $-(2(415)-9(93))=7$ \\
$2$ & $6$ & $58/13$ & $13(415)-58(93)=1$ \\
$3$ & $7$ & $415/93$ & $-(93(415)-415(93))=0$ \\
\bottomrule
\end{tabular}
\end{center}

Thus the convergent just before $415/93$ contains the certificate
\[
  1=13\cdot415-58\cdot93.
\]

\begin{center}
\begin{tikzpicture}[node distance=11mm and 8mm]
\tikzset{term/.style={draw=teal,rounded corners,fill=mist,minimum width=12mm,minimum height=8mm,font=\small},
arrow/.style={-{Stealth},navy!65,thick}}
\node[term] (a0) {$4$};
\node[term,right=of a0] (a1) {$2$};
\node[term,right=of a1] (a2) {$6$};
\node[term,right=of a2] (a3) {$7$};
\node[below=13mm of a0] (c0) {$4/1$};
\node[below=13mm of a1] (c1) {$9/2$};
\node[below=13mm of a2] (c2) {$58/13$};
\node[below=13mm of a3] (c3) {$415/93$};
\draw[arrow] (a0)--(c0); \draw[arrow] (a1)--(c1);
\draw[arrow] (a2)--(c2); \draw[arrow] (a3)--(c3);
\draw[teal!55,thick] (c0)--(c1)--(c2)--(c3);
\node[muted,below=6mm of c2] {Each convergent remembers all earlier partial quotients.};
\end{tikzpicture}
\end{center}

\begin{theorem}[Determinant identity]
Consecutive convergents satisfy
\[
p_kq_{k-1}-p_{k-1}q_k=(-1)^{k-1}.
\]
Consequently, each $p_k/q_k$ is in lowest terms and
\[
\abs{\frac{p_k}{q_k}-\frac{p_{k-1}}{q_{k-1}}}=\frac1{q_kq_{k-1}}.
\]
\end{theorem}

\begin{proof}
Take determinants in \eqref{eq:cf-matrix-product}. The determinant on the
left is $(-1)^{k+1}$, while the determinant on the right is
$p_kq_{k-1}-p_{k-1}q_k$. Since $(-1)^{k+1}=(-1)^{k-1}$, the identity follows.
\end{proof}

\section{Infinite continued fractions}

For an irrational real number $x$, repeatedly take integer parts:
\[
a_0=\floor{x},\qquad x_1=\frac1{x-a_0},\qquad a_1=\floor{x_1},\quad\dots
\]
This produces $x=[a_0;a_1,a_2,\dots]$.

\begin{example}[The silver staircase]
Let $x=\sqrt2$. Since
\[
\frac1{\sqrt2-1}=\sqrt2+1=2+(\sqrt2-1),
\]
the same remainder repeats, so
\[
\sqrt2=[1;\overline{2}]=[1;2,2,2,\dots].
\]
Its first convergents are
\[
1,\quad \frac32,\quad\frac75,\quad\frac{17}{12},\quad\frac{41}{29},\dots
\]
They alternately fall below and above $\sqrt2$.
\end{example}

\begin{theorem}[The alternating squeeze]
Let $c_k=p_k/q_k$ be the zero-indexed convergents of an irrational number
$x=[a_0;a_1,a_2,\dots]$. Then
\[
  c_0<c_2<c_4<\cdots<x<\cdots<c_5<c_3<c_1.
\]
Thus, when the convergents are counted starting with the first, the odd-numbered
ones increase and stay below $x$, while the even-numbered ones decrease and
stay above $x$.
\end{theorem}

\begin{proof}[The M\"obius derivative trick]
Fix $k$ and leave the tail variable. The continued fraction with prefix through
$a_k$ is the fractional linear, or M\"obius, map
\[
  M_k(t)=[a_0;a_1,\dots,a_k,t]
  =\frac{p_kt+p_{k-1}}{q_kt+q_{k-1}}.
\]
The determinant identity gives
\[
  M_k'(t)=
  \frac{p_kq_{k-1}-p_{k-1}q_k}{(q_kt+q_{k-1})^2}
  =\frac{(-1)^{k+1}}{(q_kt+q_{k-1})^2}.
\]
The complete tail $\alpha_{k+1}=[a_{k+1};a_{k+2},\dots]$ is finite and positive,
and
\[
  x=M_k(\alpha_{k+1}),\qquad c_k=M_k(\infty).
\]
When $k$ is even, $M_k$ is decreasing, so $c_k<x$; when $k$ is odd,
$M_k$ is increasing, so $x<c_k$. The same comparison with
$c_{k+2}=M_k([a_{k+1};a_{k+2}])$ shows that $c_k<c_{k+2}$ for even $k$ and
$c_{k+2}<c_k$ for odd $k$. These are exactly the two monotone chains displayed
above.
\end{proof}

\begin{bigidea}
Every new convergent switches sides, but every second convergent moves inward.
Continued fractions therefore trap $x$ between a rising sequence of lower
bounds and a falling sequence of upper bounds.
\end{bigidea}

\begin{center}
\begin{tikzpicture}[x=.95cm,y=1cm]
  \draw[-{Stealth},navy] (-.3,0)--(8.05,0);
  \foreach \x/\lab/\col in {0/1/teal,1.3/{7/5}/teal,2.6/{41/29}/teal,
    3.8/{\sqrt2}/gold,5.0/{99/70}/coral,6.3/{17/12}/coral,7.6/{3/2}/coral}{
    \draw[\col,thick] (\x,-.11)--(\x,.11);
    \node[\col!80!black,above=2pt] at (\x,.11) {$\lab$};
  }
  \draw[-{Stealth},teal,thick] (0,-.42)--(1.18,-.42);
  \draw[-{Stealth},teal,thick] (1.3,-.42)--(2.48,-.42);
  \draw[-{Stealth},coral,thick] (7.6,-.42)--(6.42,-.42);
  \draw[-{Stealth},coral,thick] (6.3,-.42)--(5.12,-.42);
\end{tikzpicture}

\small Schematic, not to scale: both parity subsequences move inward toward $\sqrt2$.
\end{center}

\begin{theorem}[Approximation by convergents]
If $p_k/q_k$ is a convergent of an irrational $x$, then
\[
\abs{x-\frac{p_k}{q_k}}<\frac1{q_kq_{k+1}}<\frac1{q_k^2}.
\]
Moreover, if $0<q<q_{k+1}$, then no fraction $p/q$ is closer to $x$ than both neighboring convergents in the appropriate interval. In this precise sense, convergents are best approximations with small denominators.
\end{theorem}

The determinant identity is the engine: distinct reduced fractions $a/b,c/d$ are at least $1/(bd)$ apart. There is simply not enough room for another small-denominator fraction between a convergent and $x$.

\begin{theorem}[Legendre's approximation criterion]
Let $x$ be irrational and let $p/q$ be in lowest terms with $q>0$. If
\[
  \abs{x-\frac pq}<\frac{1}{2q^2},
\]
then $p/q$ is a convergent of $x$.
\end{theorem}

The theorem is a partial converse to the approximation theorem: an
\emph{exceptionally} good rational approximation cannot have arisen by
accident. A standard proof chooses the two consecutive convergents whose
denominators surround $q$ and writes the integer vector $(p,q)$ in their
unimodular basis. If $(p,q)$ is not one of the basis vectors, the determinant
$\pm1$ forces it to be at least $1/(2q^2)$ away from $x$. This is another place
where the integral inverse of the convergent matrix does real work.

\begin{labbox}[title={Aside: continued fractions keep the calendar},breakable=false]
A calendar must use whole days. Write the tropical year as $365+\alpha$ days,
where $\alpha\approx0.2421988$. A $q$-year cycle with $p$ leap days therefore
asks for $p/q\approx\alpha$. Now
\[
 \alpha=[0;4,7,1,3,5,\dots],
 \qquad\text{with convergents}\qquad
 \frac14,\ \frac7{29},\ \frac8{33},\ \frac{31}{128},\
 \frac{163}{673},\dots
\]
The first convergent is the Julian rule: one leap day every four years. The
convergent $31/128=0.2421875$ is remarkably accurate, but a $128$-year civil
cycle is awkward.

The Gregorian rule adds a practical constraint: organize corrections by
centuries. Writing $q=100Q$ changes the approximation problem:
\[
 q=100Q:\qquad \frac pQ\approx100\alpha\approx24.21988,
 \qquad [24;4]=\frac{97}{4}.
\]
Thus the first nonintegral convergent of $100\alpha$ gives $97$ leap days in
$400$ years and the mean length
\[
  365+\frac{97}{400}=365.2425.
\]
Operationally, the rule counts $100$ multiples of $4$, removes the four
century years, and restores the one multiple of $400$: $100-4+1=97$.

Notice that $97/400$ is not a convergent of $\alpha$ itself; it emerges after
the century constraint. This is a mathematical reconstruction, not a claim
about the reformers' historical method. It shows that ``best'' depends on which
denominators an application permits.
\end{labbox}

\begin{labbox}[title={Aside: continued fractions are roughly ``base 11''}]
Khinchin's theorem says that for almost every real number,
\[
  (a_1a_2\cdots a_n)^{1/n}\longrightarrow
  K_0\approx2.685452,
\]
the universal \emph{Khinchin constant}. For efficiency, L\'evy's companion
theorem is more direct:
\[
 q_n^{1/n}\longrightarrow
 L=e^{\pi^2/(12\log2)}\approx3.275823.
\]
Since convergent error is on the scale $q_n^{-2}$, it follows that
\[
 \abs{x-p_n/q_n}^{-1/n}\longrightarrow
 L^2=e^{\pi^2/(6\log2)}\approx10.731.
\]
That is the ``base 11'' joke: one generic continued-fraction term buys about
$\log_{10}(10.731)\approx1.0306$ decimal digits. Lochs' theorem makes the
comparison exact: $N$ known decimal digits determine about
$(6\log2\log10/\pi^2)N\approx0.97027N$ continued-fraction terms.

``Almost every'' matters: a partial quotient is not literally a fixed-base
digit, and quadratic surds such as $\sqrt2$ have exceptional periodic growth.
\end{labbox}

\begin{checkpoint}[title={Experiment: $\pi$ looks typical; $e$ does not},breakable=false]
Write
\[
 G_n(x)=(a_1a_2\cdots a_n)^{1/n},
 \qquad B_n(x)=q_n^{2/n}.
\]
Khinchin and L\'evy predict, for almost every $x$, that
$G_n(x)\to2.685452\dots$ and $B_n(x)\to10.731016\dots$. For
\[
 \pi=[3;7,15,1,292,1,1,1,2,1,3,\dots],
\]
a computation gives the following rounded values:
\begin{center}
\small
\begin{tabular}{rccc}
\toprule
$n$ & $G_n(\pi)$ & $B_n(\pi)$ & $\log_{10}B_n(\pi)$ \\
\midrule
$100$     & $2.68315$ & $10.68768$ & $1.02888$ \\
$1{,}000$ & $2.66270$ & $10.54434$ & $1.02302$ \\
$10{,}000$ & $2.66320$ & $10.57847$ & $1.02442$ \\
$100{,}000$ & $2.67956$ & $10.67835$ & $1.02850$ \\
\midrule
almost-everywhere target & $2.68545$ & $10.73102$ & $1.03064$ \\
\bottomrule
\end{tabular}
\end{center}
The wobble is real: rare large partial quotients, such as the early $292$,
produce visible jumps. A published computation using $17{,}001{,}303$ terms
found $G_n(\pi)=2.686393$. This is compelling evidence, but no theorem is known
that $\pi$ satisfies either limiting law.

The number $e$ is a useful contrast, not a second generic-looking example:
\[
 e=[2;1,2,1,1,4,1,1,6,1,1,8,\dots].
\]
For $n=3m$, its first $n$ noninitial partial quotients have product $2^m m!$.
Stirling's formula therefore gives
\[
 G_{3m}(e)=(2^m m!)^{1/(3m)}\sim(2m/e)^{1/3}\longrightarrow\infty.
\]
Also $q_{3m}\geq a_1\cdots a_{3m}=2^m m!$, so $B_{3m}(e)\to\infty$ as well.
Thus $e$ is provably exceptional for both laws, while $\pi$ is numerically
persuasive but theoretically mysterious.
\end{checkpoint}

\begin{warningbox}[breakable=false]
Legendre, Lagrange, and Galois all appear in this subject. Legendre's theorem
recognizes very good rational approximations. Lagrange's theorem characterizes
eventually periodic continued fractions of quadratic irrationals; Galois's
refinement says exactly when the period begins immediately.
\end{warningbox}

\section{A parallel periodicity: digits in base \texorpdfstring{$b$}{b}}

Continued fractions are not the first place students meet repeating expansions.
Let $0<m<q$ and write $x=m/q$ in base $b\ge2$. Long division produces digits
$d_1,d_2,\dots$ and remainders $r_0,r_1,\dots$ by
\[
  r_0=m,\qquad br_n=qd_{n+1}+r_{n+1},\qquad 0\le r_{n+1}<q.
\]
There are only $q$ possible remainders. If a remainder becomes $0$, the
expansion terminates; otherwise some remainder repeats, and from that point
the deterministic long-division process repeats. Thus every rational number
has an eventually periodic base-$b$ expansion.

The converse uses only a geometric series. If a block of $\ell$ digits with
integer value $N$ repeats immediately, then
\[
  0.\overline{d_1\cdots d_\ell}_b
  =\frac{N}{b^\ell}+\frac{N}{b^{2\ell}}+\cdots
  =\frac{N}{b^\ell-1}.
\]
An initial nonrepeating block only adds a rational number with denominator a
power of $b$.

\begin{theorem}[When the repeating block begins immediately]
If the reduced fraction $m/q$ has a purely periodic base-$b$ expansion, then
$\gcd(q,b)=1$. Conversely, if $\gcd(q,b)=1$, its expansion is purely periodic.
\end{theorem}

\begin{proof}[The direction available now]
A purely periodic expansion has the form $N/(b^\ell-1)$. Its reduced
denominator divides $b^\ell-1$, which is relatively prime to $b$.
\end{proof}

The reverse implication is deliberately postponed until Chapter~6. Euler's
theorem will give
\[
  b^{\varphi(q)}\congr1\pmod q,
\]
so $q\divides b^{\varphi(q)}-1$; rewriting $m/q$ with that denominator creates
a repeating block (possibly with leading zeros). In decimal, for example,
\[
  \frac17=0.\overline{142857},\qquad
  \frac18=0.125000\dots.
\]
The first expansion is purely periodic because $\gcd(7,10)=1$; the second is
only eventually periodic.

\section{Periodic tails are fixed points}
\label{sec:periodic-fixed-points}

The matrix language answers a question that the nested notation conceals:
what kind of number can a repeating continued fraction represent?
Throughout this discussion the continued fractions are \emph{simple}:
$a_0\in\Z$ and $a_n\ge1$ for $n\ge1$.

\begin{definition}
An infinite continued fraction is \emph{eventually periodic} if its digits
repeat after a finite prefix. It is \emph{purely periodic} if repetition
starts at the first digit:
\[
 [\overline{b_0,\ldots,b_{\ell-1}}]
 =[b_0;b_1,\ldots,b_{\ell-1},b_0,b_1,\ldots].
\]
In a purely periodic simple continued fraction, every $b_i$ is positive,
including $b_0$. Rational numbers have finite simple continued fractions;
we do not call those eventually periodic here.
\end{definition}

\begin{definition}
A \emph{quadratic irrational} (or quadratic surd) is an irrational real
root of a quadratic polynomial with integer coefficients. Its other root
is its \emph{conjugate}, written $x'$. Equivalently,
\[
 x=u+v\sqrt D,\qquad x'=u-v\sqrt D
 \quad (u,v\in\Q,\ v\ne0),
\]
where $D$ is a positive nonsquare integer. Replacing $\sqrt D$ by
$-\sqrt D$ respects addition, multiplication, and division by nonzero
numbers. In particular, $(T_M(x))'=T_M(x')$ for rational $M$.
\end{definition}

\begin{lemma}[A nontrivial fixed point is quadratic]
\label{lem:mobius-fixed-point}
If an invertible rational matrix
$M=\begin{psmallmatrix}a&b\\c&d\end{psmallmatrix}$ is not a scalar
multiple of the identity and fixes an irrational number $t$ under its
linear fractional action, then $t$ is a quadratic irrational.
\end{lemma}
\begin{proof}
Multiplying $t=(at+b)/(ct+d)$ by its denominator gives
\begin{equation}\label{eq:mobius-fixed-point}
 ct^2+(d-a)t-b=0.
\end{equation}
The coefficients are not all zero: that would say $c=b=0$ and $d=a$,
making $M$ scalar. Moreover, $c$ cannot be zero, since a nonzero linear
polynomial with rational coefficients cannot vanish at an irrational
number. Thus this is a genuine quadratic equation. Clear denominators to
obtain integer coefficients.
\end{proof}

\begin{warningbox}
The qualification \emph{nontrivial} matters. The identity transformation
fixes every real number, so $t=T_I(t)$ by itself says nothing about $t$.
For a period matrix there is no such problem: its lower-left entry is
strictly positive.
\end{warningbox}

\begin{theorem}[Lagrange]
\label{thm:lagrange-periodicity}
A real number has an eventually periodic simple continued fraction if and
only if it is a quadratic irrational.
\end{theorem}

\begin{proof}[Proof that periodicity implies a quadratic equation]
First let $x=[\overline{b_0,\ldots,b_{\ell-1}}]$. Removing one complete
period leaves the same infinite tail, so
\[
 x=T_N(x),\qquad
 N=D(b_0)\cdots D(b_{\ell-1})
   =\begin{pmatrix}A&B\\C&D_0\end{pmatrix}.
\]
Because every $b_i\ge1$, we have $B,C>0$. Also
$\det N=(-1)^\ell$. The fixed-point equation is
\[
 Cx^2+(D_0-A)x-B=0.
\]
The alternating-convergent argument also applies to an infinite
positive digit sequence: it gives a limit and an error at most
$1/(q_kq_{k+1})$. That limit cannot be rational. Indeed, the
convergents are in lowest terms and $q_k\to\infty$, so they are
eventually distinct from any fixed rational $p/q$. Their distance
from $p/q$ is then at least $1/(q q_k)$, contradicting the error
bound once $q_{k+1}>q$. Thus $x$ is quadratic.

Now allow a prefix:
\[
 y=[a_0;\ldots,a_{r-1},\overline{b_0,\ldots,b_{\ell-1}}]
   =T_M(x),\qquad M=D(a_0)\cdots D(a_{r-1}).
\]
If the prefix is empty, take $M=I$. Since $T_M$ is invertible,
substituting $x=T_{M^{-1}}(y)$ into $x=T_N(x)$ gives
\begin{equation}\label{eq:prefix-conjugation}
 \boxed{\ y=T_{MNM^{-1}}(y).\ }
\end{equation}
Both $M$ and $N$ have determinant $\pm1$, so $K=MNM^{-1}$ is an
integer matrix. It is not scalar: if $K=\lambda I$, then
$N=M^{-1}KM=\lambda I$, contrary to $C>0$.

Finally, write $M=\begin{psmallmatrix}a&b\\c&d\end{psmallmatrix}$.
The number $y$ is irrational. Otherwise the inverse formula
$x=T_{M^{-1}}(y)$ would make $x$ rational. Its denominator cannot vanish
at $y=T_M(x)$: direct substitution gives
$a-cy=(ad-bc)/(cx+d)\ne0$.
The fixed-point lemma applied to $K$ therefore gives a genuine quadratic
equation for $y$ as well.
\end{proof}

\begin{bigidea}
A repeating tail is a fixed point. Attaching a finite prefix
\emph{conjugates the transformation that fixes it}. This is the missing
algebraic step between a purely periodic tail and an eventually periodic
number; it requires only matrix multiplication and inversion.
\end{bigidea}

\begin{example}[A prefix changes the equation, not its degree]
Take
\[
 x=[\overline{1,2}],\qquad
 N=D(1)D(2)=\begin{pmatrix}3&1\\2&1\end{pmatrix}.
\]
Then $2x^2-2x-1=0$, so $x=(1+\sqrt3)/2$. Prefix the digit $1$:
\[
 y=[1;\overline{1,2}]=1+\frac1x,\qquad
 M=D(1)=\begin{pmatrix}1&1\\1&0\end{pmatrix}.
\]
The conjugated matrix is
\[
 MNM^{-1}
 =\begin{pmatrix}1&1\\1&0\end{pmatrix}
  \begin{pmatrix}3&1\\2&1\end{pmatrix}
  \begin{pmatrix}0&1\\1&-1\end{pmatrix}
 =\begin{pmatrix}2&3\\1&2\end{pmatrix}.
\]
Thus $y=(2y+3)/(y+2)$, giving $y^2=3$. Since $y>0$, $y=\sqrt3$.
The same matrix will reappear as a generator of Pell solutions.
\end{example}

\section{Why every quadratic irrational repeats}
\label{sec:quadratic-finite-states}

We now prove the other direction of Lagrange's theorem. There are two
tasks: find a bounded set of integer states, and show that the continued
fraction eventually enters it. The conjugate supplies exactly the
inequalities we need.

\begin{definition}
A quadratic irrational $x$ is \emph{reduced} if
\[
 x>1\qquad\text{and}\qquad -1<x'<0.
\]
\end{definition}

\subsection*{Step 1: every quadratic irrational has a reduced tail}
Write $x_0=x$, $a_n=\floor{x_n}$, and
$x_{n+1}=1/(x_n-a_n)$. For $n\ge1$, $x_n>1$ and $a_n\ge1$.
Conjugating the recurrence gives
\[
 x_{n+1}'=\frac1{x_n'-a_n}.
\]
Consequently, once $x_n'<0$ with $n\ge1$, the next conjugate belongs to
$(-1,0)$. Every later complete quotient is reduced too.

Why must a negative conjugate appear? Invert the matrix for the first
$n$ digits and then conjugate. For $n\ge2$ this gives
\[
 x_n'=-\frac{q_{n-2}x'-p_{n-2}}{q_{n-1}x'-p_{n-1}}.
\]
The convergents $p_k/q_k$ tend to $x$, and $x'\ne x$. Hence, for all
sufficiently large $k$, the numbers
$x'-p_k/q_k$ have the same sign as $x'-x$.
The numerator and denominator in the displayed ratio therefore have
the same sign for large $n$, since the $q_k$ are positive.
It follows that $x_n'<0$ eventually, as required.
This argument also applies when the original $x$ is negative or less
than $1$.

\subsection*{Step 2: integer states evolve by an exact recurrence}
Choose an equation $Ax^2+Bx+C=0$ with integer coefficients and $A\ne0$.
After changing all three signs if necessary, the quadratic formula reads
\[
 x=\frac{-B+\sqrt D}{2A},
 \qquad D=B^2-4AC>0.
\]
Here $D$ is a nonsquare integer; it need not be squarefree. Put
$P_0=-B$ and $Q_0=2A$. Then $Q_0\ne0$ and
$Q_0\divides D-P_0^2$, because $D-P_0^2=-4AC$.
The sign change above is important: $Q_0$ may initially be negative.

Suppose $x_n=(P_n+\sqrt D)/Q_n$ with integers $P_n,Q_n$ satisfying
$Q_n\divides D-P_n^2$. Rationalizing $1/(x_n-a_n)$ yields
\begin{equation}\label{eq:quadratic-state-recurrence}
 \boxed{\quad
 P_{n+1}=a_nQ_n-P_n,\qquad
 Q_{n+1}=\frac{D-P_{n+1}^2}{Q_n}.
 \quad}
\end{equation}
Indeed,
\[
 \frac1{x_n-a_n}
 =\frac{Q_n}{\sqrt D-P_{n+1}}
 =\frac{P_{n+1}+\sqrt D}{(D-P_{n+1}^2)/Q_n}.
\]
Since $P_{n+1}\congr-P_n\pmod{Q_n}$, the new $Q_{n+1}$ is an integer.
It is nonzero because $D$ is not a square. The equality
$D-P_{n+1}^2=Q_nQ_{n+1}$ supplies the divisibility condition for the
next step. Thus the representation persists forever with the same $D$.

\subsection*{Step 3: reduced states lie in a finite box}
For a reduced complete quotient,
\[
 x_n=\frac{P_n+\sqrt D}{Q_n}>1,\qquad
 -1<\frac{P_n-\sqrt D}{Q_n}<0.
\]
Their difference is $2\sqrt D/Q_n>1$, so
$0<Q_n<2\sqrt D$. The conjugate inequalities then give
$\sqrt D-Q_n<P_n<\sqrt D$. In particular,
\begin{equation}\label{eq:reduced-state-bounds}
 \boxed{\quad |P_n|<\sqrt D,\qquad 0<Q_n<2\sqrt D.\quad}
\end{equation}
Only finitely many integer pairs satisfy these bounds.
By Step 1 all sufficiently late states lie in this box, and by Step 2
each state determines its next state, including the digit
$a_n=\floor{(P_n+\sqrt D)/Q_n}$.
Two states must repeat. Every later state, and hence every later digit,
then repeats with the same period. This completes the proof of
Theorem~\ref{thm:lagrange-periodicity}.

\begin{example}[The states for $\sqrt3$]
Starting with $(P_0,Q_0)=(0,1)$ and $D=3$ gives
\begin{center}
\begin{tabular}{ccccc}
\toprule
$n$ & $P_n$ & $Q_n$ & $x_n$ & $a_n$\\
\midrule
$0$ & $0$ & $1$ & $\sqrt3$ & $1$\\
$1$ & $1$ & $2$ & $(1+\sqrt3)/2$ & $1$\\
$2$ & $1$ & $1$ & $1+\sqrt3$ & $2$\\
$3$ & $1$ & $2$ & $(1+\sqrt3)/2$ & $1$\\
\bottomrule
\end{tabular}
\end{center}
The initial state is not reduced. The next two are reduced and form a
cycle, giving $\sqrt3=[1;\overline{1,2}]$.
\end{example}

\section{Galois: when repetition starts immediately}
\label{sec:galois}

The fixed-point equation tells us that periodic continued fractions are
quadratic. It does \emph{not} tell us which quadratic irrationals have
no initial nonrepeating segment. That distinction is Galois's theorem.

\begin{theorem}[Galois's pure-periodicity criterion]
\label{thm:galois}
The simple continued fraction of a quadratic irrational $x$ is purely
periodic if and only if $x$ is reduced:
\[
 x>1\qquad\text{and}\qquad -1<x'<0.
\]
\end{theorem}
\begin{proof}
Suppose first that $x$ is purely periodic with period length $\ell$.
The fixed-point argument shows that it is quadratic, and Step 1 above
shows that all sufficiently late complete quotients are reduced.
But $x_{k\ell}=x_0=x$ for every $k$. Taking $k$ large proves that
$x$ itself is reduced.

Conversely, suppose $x$ is reduced. All its complete quotients remain
reduced, and the finite-state proof gives $x_i=x_j$ for some $i<j$.
To conclude that the period starts at $0$, we need one more fact:
\emph{a reduced complete quotient has at most one reduced predecessor}.

Indeed, suppose the continued-fraction step sends a reduced $u$ to $z$,
and let $a=\floor u$. Then $z=1/(u-a)$, so
\[
 -\frac1{z'}=a-u'.
\]
Since $-1<u'<0$, we have $a<-1/z'<a+1$. Therefore
\begin{equation}\label{eq:reduced-predecessor}
 \boxed{\quad
 a=\floor{-1/z'},\qquad u=a+\frac1z.
 \quad}
\end{equation}
Both the previous digit and the previous reduced number are determined
by $z$. If $i>0$, equality $x_i=x_j$ therefore forces
$x_{i-1}=x_{j-1}$. Repeating the argument gives $x_0=x_{j-i}$.
The deterministic continued-fraction algorithm now repeats from its
first step, proving pure periodicity.
\end{proof}

\begin{bigidea}
Finitely many states give \emph{eventual} periodicity.
Recovering the previous reduced state removes the possible initial
segment and gives \emph{pure} periodicity. These are different steps.
For instance, the finite map $0\mapsto1$, $1\mapsto1$ has an initial
segment when started at $0$; finiteness alone does not exclude one.
\end{bigidea}

\begin{example}[Two inequalities, both necessary]
The number $\sqrt2=[1;\overline2]$ is not reduced, because
$-\sqrt2<-1$. It is eventually but not purely periodic.
In contrast,
\[
 1+\sqrt2=[\overline2],
 \qquad \frac{1+\sqrt5}{2}=[\overline1]
\]
are reduced and purely periodic.
The number $\sqrt2/2$ has conjugate in $(-1,0)$ but is less than $1$,
so it is not purely periodic either.

A displayed prefix does not necessarily represent a genuine preperiod:
\[
 [2;\overline{1,2}]=1+\sqrt3=[\overline{2,1}].
\]
Here the conjugate $1-\sqrt3$ lies in $(-1,0)$, so Galois correctly
detects pure periodicity despite the first way of writing the digits.
\end{example}

\begin{corollary}[The conjugate reverses the period]
If $x=[\overline{b_0,\ldots,b_{\ell-1}}]$, then
\[
 -\frac1{x'}=[\overline{b_{\ell-1},\ldots,b_0}].
\]
\end{corollary}
\begin{proof}
Each $D(b_i)$ is symmetric, so reversing the digits transposes their
product $N=\begin{psmallmatrix}A&B\\C&D_0\end{psmallmatrix}$.
Its fixed-point equation is $Cx^2+(D_0-A)x-B=0$. Substituting
$x'=-1/z$ gives
\[
 Bz^2+(D_0-A)z-C=0,
\]
which is the fixed-point equation for $N^T$. Since $B,C>0$, this
quadratic has exactly one positive root. Both $-1/x'$ and the
reversed periodic continued fraction are that root.
\end{proof}

\begin{checkpoint}[title={Three names, three different conclusions}]
\textbf{Legendre}: an exceptionally close rational approximation must
be a convergent.
\quad\textbf{Lagrange}: the numbers with eventually periodic simple
continued fractions are exactly the quadratic irrationals.
\quad\textbf{Galois}: among them, pure periodicity means exactly
$x>1$ and $-1<x'<0$.
The first direction of Lagrange is the fixed-point calculation;
Galois's additional content is the precise conjugate criterion.
\end{checkpoint}


\section{Periodicity meets Pell's equation}
\label{sec:pell}

Fix a nonsquare integer $D\ge2$. \emph{Pell's equation} asks for positive
integers $p,q$ such that
\[
 p^2-Dq^2=1.
\]
We will also study the equation with right-hand side $-1$. There are two
questions to settle: why must a solution appear among the convergents of
$\sqrt D$, and which convergents actually work? An approximation estimate
answers the first; an exact identity answers the second.

\subsection*{Every Pell solution is a convergent}
Suppose $p^2-Dq^2=\pm1$, with $p,q>0$. Any common divisor of $p,q$ would
divide $1$, so $p/q$ is reduced. Factoring a difference of squares gives
\begin{equation}\label{eq:pell-error}
 \abs{\sqrt D-\frac pq}
 =\frac{1}{q^2(\sqrt D+p/q)}
 <\frac{1}{2q^2}.
\end{equation}
For the strict inequality, note that
$(p/q)^2\ge D-1\ge1$, while $\sqrt D>1$.
Legendre's criterion now makes $p/q$ a convergent of $\sqrt D$.
This does not yet prove that a solution exists: we still need a convergent
whose difference of squares is exactly $1$ or $-1$.

\subsection*{The palindrome in a square-root period}
Write $a_0=\floor{\sqrt D}$. The first complete quotient is
\[
 x_1=\frac1{\sqrt D-a_0}>1,
 \qquad x_1'=-\frac1{\sqrt D+a_0}\in(-1,0).
\]
Galois's theorem therefore makes $x_1$ purely periodic. Let $L$ be its
\emph{least} period length, so that
$\sqrt D=[a_0;\overline{a_1,\ldots,a_L}]$.

\begin{theorem}[The square-root palindrome]
\label{thm:sqrt-palindrome}
The last digit of the least period is $a_L=2a_0$, and
\[
 a_j=a_{L-j}\qquad(1\le j\le L-1).
\]
Thus
\[
 \boxed{\quad
 \sqrt D=[a_0;\overline{a_1,\ldots,a_{L-1},2a_0}],
 \quad (a_1,\ldots,a_{L-1})\text{ is a palindrome}.
 \quad}
\]
For $L=1$ the palindrome is the empty block.
\end{theorem}
\begin{proof}
Apply Corollary~\ref{cor:period-reversal} to $x_1$:
\[
 -\frac1{x_1'}=\sqrt D+a_0
   =[\overline{a_L,a_{L-1},\ldots,a_1}].
\]
Its integer part is $2a_0$, so $a_L=2a_0$. On the other hand,
\[
 \sqrt D+a_0=2a_0+\frac1{x_1}
   =[2a_0;\overline{a_1,\ldots,a_L}].
\]
These are two expansions of the same irrational number. The
continued-fraction algorithm determines its digits uniquely. Comparing
the next $L-1$ digits gives
$(a_1,\ldots,a_{L-1})=(a_{L-1},\ldots,a_1)$.
\end{proof}

For example,
\[
 \sqrt{23}=[4;\overline{1,3,1,8}],\qquad
 \sqrt{14}=[3;\overline{1,2,1,6}],\qquad
 \sqrt2=[1;\overline2].
\]
The final $2a_0$ marks the end of the symmetric block; it is not itself
part of the palindrome asserted by the theorem.

There is a useful consequence for the integer states
$x_n=(P_n+\sqrt D)/Q_n$ of
\eqref{eq:quadratic-state-recurrence}, starting with $(P_0,Q_0)=(0,1)$.
Since $x_{L+1}=x_1$,
\[
 x_L=a_L+\frac1{x_{L+1}}=2a_0+(\sqrt D-a_0)=\sqrt D+a_0.
\]
Consequently $(P_L,Q_L)=(a_0,1)$, and the same state occurs at every
positive multiple of $L$. Conversely, if $n\ge1$ and $Q_n=1$, reducedness
gives $-1<P_n-\sqrt D<0$, forcing $P_n=a_0$. Thus $x_n=x_L$.
The $L$ states in a least cycle are distinct, so
\begin{equation}\label{eq:pell-state-return}
 \boxed{\qquad Q_n=1\quad\Longleftrightarrow\quad L\divides n
 \qquad(n\ge1).\qquad}
\end{equation}

\subsection*{An exact identity for every convergent}
\begin{lemma}[The complete quotient measures the Pell remainder]
\label{lem:pell-remainder}
For every $n\ge1$,
\begin{equation}\label{eq:pell-remainder}
 \boxed{\quad p_{n-1}^2-Dq_{n-1}^2=(-1)^nQ_n.\quad}
\end{equation}
\end{lemma}
\begin{proof}
Keep the entire tail in the matrix product:
\[
 \sqrt D
 =\frac{p_{n-1}x_n+p_{n-2}}{q_{n-1}x_n+q_{n-2}},
 \qquad x_n=\frac{P_n+\sqrt D}{Q_n}.
\]
Substitute the second expression into the first and multiply through.
Because $\sqrt D$ is irrational, the rational coefficients and the
coefficients of $\sqrt D$ must agree separately. They give
\begin{align*}
 p_{n-1}&=P_nq_{n-1}+Q_nq_{n-2},\\
 Dq_{n-1}&=P_np_{n-1}+Q_np_{n-2}.
\end{align*}
Multiply the first equality by $p_{n-1}$ and the second by $q_{n-1}$,
then subtract:
\begin{align*}
 p_{n-1}^2-Dq_{n-1}^2
 &=Q_n(p_{n-1}q_{n-2}-p_{n-2}q_{n-1})\\
 &=(-1)^nQ_n.
\end{align*}
The last step is the determinant of the product of the $n$ digit
matrices $D(a_0),\ldots,D(a_{n-1})$. The initial conventions
$p_{-1}=1,q_{-1}=0$ also cover $n=1$.
\end{proof}

\begin{theorem}[Exactly which convergents solve Pell]
\label{thm:pell-convergents}
For each $m\ge1$,
\begin{equation}\label{eq:pell-period}
 p_{mL-1}^2-Dq_{mL-1}^2=(-1)^{mL}.
\end{equation}
These are all positive integer solutions to the two equations
$p^2-Dq^2=\pm1$. In particular:
\begin{itemize}
 \item If $L$ is even, $+1$ occurs at every index $mL-1$, and the
 equation with right-hand side $-1$ has no integer solution.
 \item If $L$ is odd, $-1$ occurs for odd $m$ and $+1$ for even $m$.
\end{itemize}
The least positive solution of the $+1$ equation is
$(p_{L-1},q_{L-1})$ for even $L$ and $(p_{2L-1},q_{2L-1})$ for odd $L$.
\end{theorem}
\begin{proof}
All $Q_n$ with $n\ge1$ are positive. The lemma says that a convergent
solves one of the two equations exactly when $Q_n=1$.
By \eqref{eq:pell-state-return}, this happens exactly when $n=mL$;
\eqref{eq:pell-period} and its sign follow. Legendre's criterion, already
applied above, ensures that no positive solution lies outside the
convergents. Finally, the recurrence makes $p_0,p_1,\ldots$ strictly
increasing, so the first admissible index gives the least $p$, and hence
the least $q$ as well. If a solution of the $-1$ equation had either
coordinate zero, it would require $p^2=-1$ or $Dq^2=1$, both impossible;
changing signs reduces every other integer solution to a positive one.
\end{proof}

\begin{warningbox}[title={Stop before the final digit}]
The first candidate is
\[
 \frac{p_{L-1}}{q_{L-1}}=[a_0;a_1,\ldots,a_{L-1}],
\]
which stops \emph{before} the final $2a_0$. There are $L$ digits here,
counting $a_0$. For $\sqrt3=[1;\overline{1,2}]$, the candidate is
$[1;1]=2$, and $2^2-3\cdot1^2=1$. Including the next digit gives
$[1;1,2]=5/3$, whose remainder is $25-27=-2$.
\end{warningbox}

\begin{example}[An odd period: $D=13$]
The state calculation gives
$\sqrt{13}=[3;\overline{1,1,1,1,6}]$, with $L=5$.
The lemma explains every entry of the last column below.
\begin{center}
\begin{tabular}{ccrrr}
\toprule
$k$ & $a_k$ & $p_k/q_k$ & $Q_{k+1}$ & $p_k^2-13q_k^2$\\
\midrule
0 & 3 & $3/1$ & 4 & $-4$\\
1 & 1 & $4/1$ & 3 & $3$\\
2 & 1 & $7/2$ & 3 & $-3$\\
3 & 1 & $11/3$ & 4 & $4$\\
4 & 1 & $18/5$ & 1 & $-1$\\
5 & 6 & $119/33$ & 4 & $4$\\
6 & 1 & $137/38$ & 3 & $-3$\\
7 & 1 & $256/71$ & 3 & $3$\\
8 & 1 & $393/109$ & 4 & $-4$\\
9 & 1 & $649/180$ & 1 & $1$\\
\bottomrule
\end{tabular}
\end{center}
Thus the first negative solution is $(18,5)$, while the first positive
one is $(649,180)$. Squaring the former gives the latter:
\[
 (18+5\sqrt{13})^2=649+180\sqrt{13}.
\]
\end{example}

\subsection*{The period matrix becomes a Pell matrix}
Set $p=p_{L-1}$, $q=q_{L-1}$ and
$M=D(a_0)\cdots D(a_{L-1})$. At the returning state $P_L=a_0,Q_L=1$,
the two coefficient identities in the lemma give
\[
 M\begin{pmatrix}1&a_0\\0&1\end{pmatrix}
 =\begin{pmatrix}p&a_0p+p_{L-2}\\q&a_0q+q_{L-2}\end{pmatrix}
 =\begin{pmatrix}p&Dq\\q&p\end{pmatrix}.
\]
The shift by $a_0$ accounts for the fact that $M$ sends the tail
$\sqrt D+a_0$ to $\sqrt D$. The resulting matrix fixes $\sqrt D$ and
has determinant $p^2-Dq^2=(-1)^L$. For odd $L$, square it to obtain
determinant $+1$.

\subsection*{Why one fundamental solution generates all the others}
For $z=r+s\sqrt D$, define its \emph{norm} by
$N(z)=zz'=r^2-Ds^2$. Conjugation respects multiplication, so
$N(zw)=N(z)N(w)$. If $(u,v)$ is the least positive solution of the $+1$
equation, put $\varepsilon=u+v\sqrt D>1$. Its inverse is
$u-v\sqrt D$, again with integer coefficients.

\begin{proposition}[All positive Pell solutions are powers]
\label{prop:pell-powers}
Every positive solution of $p^2-Dq^2=1$ has the form
\[
 p+q\sqrt D=\varepsilon^j\qquad(j=1,2,\ldots).
\]
\end{proposition}
\begin{proof}
Powers have integer coefficients, norm $1$, and positive coefficients
for $j\ge1$. Conversely, let $\alpha=p+q\sqrt D>1$ have norm $1$.
Choose $j\ge0$ with $\varepsilon^j\le\alpha<\varepsilon^{j+1}$ and put
$\beta=\alpha\varepsilon^{-j}=r+s\sqrt D$. Then $r,s$ are integers,
$N(\beta)=1$, and $1\le\beta<\varepsilon$. Since $\beta'=1/\beta$,
\[
 r=\frac{\beta+\beta^{-1}}2>0,
 \qquad s=\frac{\beta-\beta^{-1}}{2\sqrt D}\ge0.
\]
If $\beta>1$, the function $t-t^{-1}$ is strictly increasing for $t>0$,
so $0<s<v$. This would be a smaller positive Pell solution, a
contradiction. Therefore $\beta=1$ and $\alpha=\varepsilon^j$.
Since $q>0$, we have $j\ge1$.
\end{proof}

Equivalently, multiplication by $\varepsilon$ is the matrix recurrence
\[
 \begin{pmatrix}p^{(j+1)}\\q^{(j+1)}\end{pmatrix}
 =\begin{pmatrix}u&Dv\\v&u\end{pmatrix}
  \begin{pmatrix}p^{(j)}\\q^{(j)}\end{pmatrix},
 \qquad (p^{(0)},q^{(0)})=(1,0).
\]
The superscript counts powers; the subscript on $p_k/q_k$ counts
continued-fraction convergents. For $D=2$, $\varepsilon=3+2\sqrt2$ gives
$(3,2),(17,12),(99,70)$.

\begin{bigidea}
Periodicity gives the exact return $Q_{mL}=1$. The determinant identity
then turns that return into $p_{mL-1}^2-Dq_{mL-1}^2=(-1)^{mL}$.
Legendre proves completeness, and powers of the least positive solution
give a second way to list every positive solution. Each step has a
different job.
\end{bigidea}

\begin{checkpoint}
Compute $99/70$ as a decimal. Its accuracy follows from
$99^2-2\cdot70^2=1$, which forces
\[
 \abs{\sqrt2-99/70}=\frac1{70^2(\sqrt2+99/70)}.
\]
For a second check, use the period of $\sqrt{23}$ above to find the
fundamental solution $(24,5)$ without searching through denominators.
\end{checkpoint}

\section{The Cattle of the Sun: Archimedes' challenge}
\label{sec:archimedes-cattle}

The cattle problem attributed to Archimedes asks for a herd belonging to
the Sun, with bulls and cows in four colors: white, black, yellow, and
dappled. Its proportions can be written as linear equations. Two further
conditions ask for a square and a triangle. Together they lead to a Pell
equation with a remarkably large first solution.

This optional application begins with the
\href{https://math.nyu.edu/Archimedes/Cattle/Statement.html}{ancient problem's conditions}.
For a detailed modern exposition, see Ilan Vardi,
\href{https://doi.org/10.1080/00029890.1998.12004887}{``Archimedes' Cattle Problem,''
\emph{American Mathematical Monthly} \textbf{105} (1998), 305--319};
an \href{https://www.maths.tcd.ie/~lebed/Cattle\%20problem\%20Ilan\%20Vardi.pdf}{accessible preprint}
is also available. The calculations below can be reproduced with exact
integer arithmetic in the first notebook.

\subsection*{Eight unknowns, seven linear conditions}
Let $W,B,Y,S$ count white, black, yellow, and spotted (dappled) bulls;
let $w,b,y,s$ count cows of the corresponding colors. The conditions are
\begin{align*}
 W&=\left(\frac12+\frac13\right)B+Y,
 &B&=\left(\frac14+\frac15\right)S+Y,
 &S&=\left(\frac16+\frac17\right)W+Y,\\
 w&=\left(\frac13+\frac14\right)(B+b),
 &b&=\left(\frac14+\frac15\right)(S+s),\\
 y&=\left(\frac16+\frac17\right)(W+w),
 &s&=\left(\frac15+\frac16\right)(Y+y).
\end{align*}
The cow equations use an entire herd of the specified color, including
its bulls. The two additional conditions concern bulls only:
\begin{equation}\label{eq:cattle-shapes}
 W+B=\text{a square},\qquad
 Y+S=\frac{h(h+1)}2\quad\text{for some integer }h\ge1.
\end{equation}

Eliminate $B,Y,S$ from the first three equations to obtain
\[
 B=\frac{267}{371}W,\qquad
 Y=\frac{297}{742}W,\qquad
 S=\frac{790}{1113}W.
\]
The four cow equations form a cycle. Substituting around that cycle
expresses $w$ as a known multiple of $W$ plus
$\frac7{12}\frac9{20}\frac{11}{30}\frac{13}{42}w$.
Since this coefficient is less than $1$, $w$, and then every cow count,
is uniquely determined by $W$. Clearing denominators gives
\begin{center}
\begin{tabular}{lrr}
\toprule
Color & Bulls & Cows\\
\midrule
White   & $10\,366\,482t$ & $7\,206\,360t$\\
Black   & $7\,460\,514t$  & $4\,893\,246t$\\
Yellow  & $4\,149\,387t$  & $5\,439\,213t$\\
Dappled & $7\,358\,060t$  & $3\,515\,820t$\\
\bottomrule
\end{tabular}
\end{center}
Substitution verifies all seven equations, and the eight coefficients
have gcd $1$. Therefore all eight counts are integers exactly when
$t$ is an integer: a B\'ezout combination of those counts equals $t$.
Every positive solution of the linear part has $t\ge1$, and the total
number of cattle is
\begin{equation}\label{eq:cattle-total-scale}
 T=50\,389\,082t.
\end{equation}

\subsection*{The square and the triangle produce Pell}
The square condition is
\[
 W+B=17\,826\,996t
 =4\cdot4\,456\,749t,
 \qquad4\,456\,749=3\cdot11\cdot29\cdot4657.
\]
The four factors in the last product are distinct primes (trial division
checks $4657$). Unique factorization therefore says that this is a square
exactly when
\[
 t=4\,456\,749r^2\qquad(r\ge1).
\]
Indeed, every prime in the displayed squarefree product must have odd
exponent in $t$, and every other prime must have even exponent.
Now the triangular condition becomes
\[
 11\,507\,447t=\frac{h(h+1)}2,
 \qquad
 51\,285\,802\,909\,803r^2=\frac{h(h+1)}2.
\]
Multiplying by $8$ and completing the square gives
\begin{equation}\label{eq:cattle-large-pell}
 \boxed{\quad X^2-410\,286\,423\,278\,424r^2=1,
 \qquad X=2h+1.\quad}
\end{equation}
Conversely, a positive solution of this equation has odd $X$, since the
coefficient of $r^2$ is divisible by $8$. Thus $h=(X-1)/2$ is an integer
and reconstructs the triangle. Together with the table, this proves an
equivalence, not just a necessary condition.

The coefficient is large, but contains a large square factor:
\[
 410\,286\,423\,278\,424=4\,729\,494\cdot9314^2.
\]
Put $Z=9314r$. The equivalent smaller-coefficient problem is
\begin{equation}\label{eq:cattle-small-pell}
 \boxed{\quad X^2-4\,729\,494Z^2=1,
 \qquad 9314\divides Z.\quad}
\end{equation}
Keeping the divisibility condition is essential: an arbitrary solution
of the smaller Pell equation need not give an integer $r$.

\subsection*{A short period, followed by a divisibility test}
The integer-state algorithm gives $a_0=2174$ and least period length
$L=92$ for $\sqrt{4\,729\,494}$. The final digit is $4348$; the preceding
$91$ digits form a palindrome. Since $92$ is even, the convergent
$p_{91}/q_{91}$ gives the fundamental solution
\begin{align*}
 u&=109931986732829734979866232821433543901088049,\\
 v&=50549485234315033074477819735540408986340.
\end{align*}
These values can be checked by the exact equality
$u^2-4\,729\,494v^2=1$. Write
$\varepsilon=u+v\sqrt{4\,729\,494}$ and
\[
 \varepsilon^j=X_j+Z_j\sqrt{4\,729\,494}.
\]
Proposition~\ref{prop:pell-powers} proves that this lists every positive
solution of the smaller Pell equation. We need the first $j$ for which
$9314\divides Z_j$.

We can find it without ever forming the enormous powers. Reduce the
multiplication recurrence modulo $9314$. Here
\[
 u\congr9063,\quad v\congr7708,\quad
 4\,729\,494\congr7296\pmod{9314}.
\]
The following finite computation checks every exponent up to the first
success, using only small integers:
\begin{quote}\begin{minipage}{\linewidth}\small
\begin{verbatim}
modulus = 9314
u, v, d = 9063, 7708, 7296
a, b = 1, 0
for j in range(1, 2330):
    a, b = ((u*a + d*v*b) % modulus,
            (v*a + u*b) % modulus)
    if b == 0:
        print(j, a, b)
        break
\end{verbatim}
\end{minipage}\end{quote}
Its output is \texttt{2329 9313 0}; no earlier exponent passes the test.
Thus $\varepsilon^{2329}$ produces the smallest admissible $r$ and
hence the smallest herd. This is a reproducible finite calculation,
distinct from the general Pell proofs above. Vardi explains a shorter
divisibility test using a quadratic finite field.

\subsection*{An enormous answer with a compact description}
Let $X+Z\sqrt{4\,729\,494}=\varepsilon^{2329}$. Then the least total is
\begin{equation}\label{eq:cattle-exact-total}
 T=50\,389\,082\cdot4\,456\,749
       \left(\frac{Z}{9314}\right)^2.
\end{equation}
This is an exact specification using integer arithmetic. The first
notebook computes $X,Z$ by repeated squaring, verifies the seven
proportions and both shape conditions, and checks
\[
 10^{206544}\le T<10^{206545}.
\]
The decimal expansion of the smallest herd's size therefore has
$206\,545$ digits:
\[
 T\approx7.760271406\times10^{206544}.
\]
Its size is a mathematical consequence of the conditions, rather than
a plausible census of animals.

Conjugation also gives a compact real-number expression, of the kind
developed in Vardi's exposition. Since
\[
 Z=\frac{\varepsilon^{2329}-\varepsilon^{-2329}}
          {2\sqrt{4\,729\,494}},
\]
substitution into \eqref{eq:cattle-exact-total} yields
\[
 T=\frac{25\,194\,541}{184\,119\,152}
       \left(\varepsilon^{4658}+\varepsilon^{-4658}-2\right).
\]
An expression can be short even when its decimal expansion is long.

\begin{bigidea}
The linear conditions leave one scale factor. Prime-exponent parity
makes that factor a fixed number times a square. The triangular
condition becomes Pell, and a divisibility test selects the allowed
powers of its fundamental solution. The same elementary tools that
solve $p^2-2q^2=1$ organize the cattle problem.
\end{bigidea}


\begin{exercisebox}
\begin{enumerate}
  \item Find the continued fraction of $355/113$. List all convergents. \exlevel{Fluency}
  \item Compute the first six convergents of $[1;\overline{1}]$. What familiar sequence appears? \exlevel{Fluency}
  \item Prove the determinant identity by induction. Use it to prove $\gcd(p_k,q_k)=1$. \exlevel{Reasoning}
  \item Show that $[a_0;\dots,a_n]=[a_0;\dots,a_n-1,1]$ when $a_n>1$. Explain the only ambiguity in finite expansions. \exlevel{Reasoning}
  \item Derive $\sqrt3=[1;\overline{1,2}]$ and compute four convergents. \exlevel{Reasoning}
  \item Write the tropical year as $365+\alpha$ with
  $\alpha\approx0.2421988$. Compute the first four nonzero convergents of
  $\alpha$ and identify the Julian rule. Then impose a century-based cycle by
  approximating $100\alpha$; explain how $97/4$ becomes the Gregorian rule
  $97/400$. Compare it with $31/128$. \exlevel{Exploration}
  \item Investigate which convergents to $\sqrt2$ solve $p^2-2q^2=1$ and which solve $-1$. State and prove a pattern using the recurrence. \exlevel{Challenge}
  \item Let $A/B=[a_0;\dots,a_n]$ and use the convergent recurrences to prove
  $r_k=(-1)^k(q_kA-p_kB)$. Apply the penultimate convergent of $527/341$ to
  recover a B\'ezout identity for $\gcd(527,341)$. \exlevel{Reasoning}
  \item For $x=\sqrt3$, compute enough convergents to display three lower
  bounds and three upper bounds. Verify that the lower bounds increase and the
  upper bounds decrease, then explain why the M\"obius derivative has the sign
  needed for each chain. \exlevel{Reasoning}
  \item Prove the matrix product formula \eqref{eq:cf-matrix-product} by
  induction. Explain why its inverse always has integer entries. \exlevel{Reasoning}
  \item Long-divide $1/7$ and $5/12$ in base $10$, recording the remainder at
  every step. Use repeated remainders to predict the preperiod and period.
  \exlevel{Fluency}
  \item Prove directly from a geometric series that every eventually periodic
  base-$b$ expansion is rational. \exlevel{Reasoning}
  \item After studying Euler's theorem, return to prove that $\gcd(q,b)=1$
  implies the base-$b$ expansion of $m/q$ is purely periodic. Give an explicit
  upper bound for a period length. \exlevel{Reasoning}
  \item Verify that $355/113$ satisfies Legendre's inequality for $\pi$ using
  a calculator value of $\pi$. What conclusion does the theorem certify?
  \exlevel{Exploration}
  \item Let $t=[\overline{1,2}]$. Use the matrix $D(1)D(2)$ to find a quadratic
  equation for $t$, and identify the positive root. \exlevel{Reasoning}
  \item Check Galois's inequalities for $1+\sqrt2$ and
  $(1+\sqrt5)/2$. Explain why their continued fractions are purely periodic.
  \exlevel{Reasoning}
  \item Let $p^2-Dq^2=\pm1$ with $p,q>0$. Derive
  \eqref{eq:pell-error}, verify its strict Legendre bound, and conclude that
  $p/q$ is a convergent of $\sqrt D$. \exlevel{Reasoning}
  \item Compute the period of $\sqrt3$. Use its length to predict whether one
  period solves the $+1$ or $-1$ Pell equation. Then use powers of
  $\begin{psmallmatrix}2&3\\1&2\end{psmallmatrix}$ to generate three solutions
  of $p^2-3q^2=1$. \exlevel{Exploration}
  \item Define $G_n(x)=(a_1\cdots a_n)^{1/n}$ and $B_n(x)=q_n^{2/n}$.
  Use the Sage notebook to reproduce several rows of the $\pi$ table above.
  Explain why agreement with the Khinchin--L\'evy constants is evidence but
  not proof. Then derive the displayed formula for $G_{3m}(e)$ directly from
  the pattern in the continued fraction of $e$. \exlevel{Exploration}
  \item Compute $D(1)D(2)$ and $D(2)D(1)$, and write the linear
  fractional transformation associated to each. Explain the order of
  composition and find each transformation's positive fixed point.
  \exlevel{Fluency}
  \item Let $x=[\overline{1,2}]$ and $N=D(1)D(2)$. For $j=1,2$, set
  $M=D(j)$ and compute $MNM^{-1}$. Use its fixed-point equation to
  identify $y=[j;\overline{1,2}]$. Which of these numbers is purely
  periodic? Write its period without an initial prefix.
  \exlevel{Reasoning}
  \item Prove that an invertible rational linear fractional transformation
  sends every irrational real number to an irrational real number.
  Explain why a pole cannot occur at an irrational input. Then identify
  the missing hypothesis in the assertion that $t=T_M(t)$ always makes
  $t$ quadratic. Give a counterexample. \exlevel{Reasoning}
  \item Use Galois's criterion to decide which of
  \[
   \sqrt2,\quad 1+\sqrt2,\quad \frac{1+\sqrt5}{2},\quad
   \frac{1+\sqrt3}{2},\quad \frac{\sqrt2}{2},\quad -\sqrt2
  \]
  have purely periodic continued fractions. Explain why all six are
  eventually periodic. \exlevel{Fluency}
  \item Starting from $D=13$ and $(P_0,Q_0)=(0,1)$, use
  \eqref{eq:quadratic-state-recurrence} to compute the continued fraction
  of $\sqrt{13}$. Record $(P_n,Q_n,a_n)$ until a state first repeats,
  and verify the reduced-state bounds after the first step.
  \exlevel{Exploration}
  \item Recover the unique reduced predecessor of
  $z=(1+\sqrt3)/2$ using \eqref{eq:reduced-predecessor}.
  For $D=3$, enumerate the integer pairs satisfying
  \eqref{eq:reduced-state-bounds} that actually represent reduced surds
  and satisfy $Q\divides3-P^2$. Show how the continued-fraction step
  permutes them. Why would finiteness without reversibility not suffice
  to prove pure periodicity? \exlevel{Reasoning}
  \item Let $x=[\overline{1,2,3}]$. Compute its period matrix and its
  positive fixed point. Transpose the matrix to obtain the equation for
  $[\overline{3,2,1}]$, and verify directly that this number is $-1/x'$.
  \exlevel{Reasoning}
  \item Compute the integer states and least period for $\sqrt{23}$.
  Identify the palindrome, find the fundamental Pell solution, and
  calculate what happens if you include the final digit $2a_0$ in the
  first candidate. \exlevel{Fluency}
  \item For $D=6$, calculate four convergents and the corresponding
  $Q_{k+1}$. Verify \eqref{eq:pell-remainder} in each row. Use
  \eqref{eq:pell-state-return} to classify every convergent that solves
  either sign of Pell. \exlevel{Reasoning}
  \item Decide whether $p^2-Dq^2=-1$ is solvable for
  $D=2,3,5,7,13,14$. Give the least positive solution when it exists.
  For $D=3$, also give a one-line impossibility proof modulo $3$.
  \exlevel{Reasoning}
  \item Let $p^2-3q^2=1$ with $p,q>0$. Show that multiplying
  $p+q\sqrt3$ by $2-\sqrt3$ gives $(1,0)$ from $(2,1)$,
  and otherwise a positive solution with smaller second coordinate.
  Use this descent to prove directly that all positive solutions are
  powers of $2+\sqrt3$. \exlevel{Reasoning}
  \item Use the integer-state algorithm to find the least period of
  $\sqrt{61}$. Find the first negative and positive Pell solutions,
  and explain why the theorem proves that each is the first of its
  sign. Use exact integers, not a decimal approximation to $\sqrt{61}$.
  \exlevel{Exploration}
  \item Verify the cattle table at $t=1$, including its gcd and total.
  Explain why the seven linear equations have only one free scale
  factor, and why the gcd certifies that this scale is an integer.
  Does the table at $t=1$ solve the full problem? \exlevel{Reasoning}
  \item For positive integers $c,r$, prove that $cr^2$ is triangular
  exactly when $X^2-8cr^2=1$ has a positive odd solution $X$.
  Apply this to the cattle problem, and prove both directions of the
  passage from \eqref{eq:cattle-large-pell} to
  \eqref{eq:cattle-small-pell}. \exlevel{Reasoning}
  \item Reproduce the modular search for the cattle problem in the
  first notebook. Explain why its first hit proves minimality. Using
  $\varepsilon^{2329}\congr-1\pmod{9314}$ and the absence of an
  earlier hit, prove that the admissible exponents are exactly the
  positive multiples of $2329$. Why does $4658$ appear in the total's
  compact formula? \exlevel{Exploration}
\end{enumerate}
\end{exercisebox}
